\documentclass[../../main.tex]{subfiles}
\begin{document}

\begin{center}
  {\fontsize{48pt}{58pt}\selectfont \textbf{1972}\normalsize\textbf{年度}}
\end{center}

\begin{table}[htb]
  \centering
  \caption{出題テーマと難易度}
  \begin{tabular}{|c|p{9cm}|c|}
    \hline
    \textbf{問題番号} & \textbf{テーマ} & \textbf{難易度} \\
    \hline
    第1問 & ベクトル加法の交換・結合法則と帰納法 & \\
    \hline
    第2問 & 有理関数の定積分と対数の極限 & \\
    \hline
    第3問 & 三数の和と三乗和から得る等式 & \\
    \hline
    第4問 & 三角形内の平行四辺形とベクトルの係数 & \\
    \hline
    第5問 & 楕円上の動点が掃く面積と速度 & \\
    \hline
    第6問 & 2枚のカードの和の期待値 & \\
    \hline
  \end{tabular}
\end{table}

\begin{table}[htb]
  \centering
  \caption{解答}
  \renewcommand{\arraystretch}{1.4}
  \begin{tabular}{|c|c|p{9cm}|}
    \hline
    \textbf{問題番号} & \textbf{小問} & \textbf{答え} \\
    \hline
    第1問 & -- & 順序によらずベクトルの和は等しい（証明問題） \\
    \hline
    第2問 & -- & $-\dfrac{3}{2}\log 3$ \\
    \hline
    第3問 & -- & $x,y,z$の少なくとも1つは$a$に等しい（証明問題） \\
    \hline
    第4問 & -- &  \\
    \hline
    第5問 & -- & $u=\dfrac{2t}{ab}$ \\
    \hline
    第6問 & -- & $F=2E$（証明問題） \\
    \hline
  \end{tabular}
\end{table}

\paragraph{試験形態}
6問．試験時間は要確認．

\end{document}
